\begin{titlepage}
    \begin{center}
        \vspace{-2cm}
        \centerline{\includegraphics[height=4cm]{logo}}
        \vspace{-0.5cm}

        {\nazzz \textbf{دانشگاه اصفهان}}
        \vspace{-.3cm}

        {\nazz \textbf{دانشکده ریاضی و آمار}}
        \vspace{-.3cm}

        {\naz \textbf{گروه علوم کامپیوتر}}

        \vspace{1.2cm}
        \textbf{\nazzzzz گزارش پیشرفت پروژه کارشناسی رشته علوم کامپیوتر}

        \vspace{1cm}
        \textbf{\Titr مطالعه و ارزیابی روش‌های بهینه‌سازی \lr{k-means}\\
            با تمرکز بر الگوریتم \lr{GLKM}}

        \vspace{1.2cm}
        \textbf{\nazzzz{استاد راهنما:}}
        \vspace{.2cm}

        \textit{\nazzz {سرکار خانم دکتر نجمه حسینی منجزی}}

        \vspace{1cm}
        \textbf{\nazzzz{دانشجو:}}
        \vspace{.2cm}

        \textit{\nazzz {داوود نصرتی امیرآبادی}}

        \vspace{.5cm}
        \textbf{\nazzz{مرداد ۱۴۰۵}}
    \end{center}
\end{titlepage}

\chapter*{چکیده}
این گزارش، فعالیت‌های انجام‌شده در مرحله‌ی نخست پروژه‌ی کارشناسی با موضوع
خوشه‌بندی \lr{k-means} و الگوریتم \lr{Greedy Local K-Means (GLKM)} را
مستند می‌نماید. ابتدا سیر تاریخی و الگوریتمی چند روش اصلی، از الگوریتم
مک‌کوئین و کوانتیزاسیون لوید تا روش هارتیگان--وونگ و مقداردهی اولیه‌ی
\lr{k-means++}، مطالعه و مقایسه شده است. سپس ساختار \lr{GLKM} و دو ایده‌ی
توسعه‌ای بررسی شده‌اند: برآورد خودکار تعداد خوشه‌ها با نمودار \lr{elbow}، و
استفاده از یک افراز اولیه‌ی حاصل از \lr{k-means++} پیش از اجرای \lr{GLKM}.
برای ایده‌ی نخست، ابزار اجرای خودکار روی چند مقدار $k$ و تخمین نقطه‌ی آرنج
پیاده‌سازی شد. آزمایش‌ها نشان دادند که در
داده‌ی \lr{Gisette} منحنی هزینه تقریباً خطی بود و تخمین elbow با ساختار
برچسب‌های واقعی سازگار نبود. در ایده‌ی دوم نیز،
کیفیت خوشه‌بندی در آزمایش انجام‌شده کاهش یافت. بنابراین این مرحله
هنوز به یک بهبود عمومی برای \lr{GLKM} منجر نشده است، اما عملکرد دو
مسیر پیشنهادی را روشن و بستر آزمایش‌های بعدی را آماده کرده است.

\vspace{0.75cm}
\textbf{کلیدواژه‌ها}:
خوشه‌بندی، \lr{k-means}، \lr{GLKM}، روش \lr{elbow}، \lr{k-means++}
